\chapter{来源与使用边界}
\label{app:sources}

\section{怎样使用本卷的来源}

本卷不把材料排成一张固定的可靠性阶梯。一个诏令最适合说明朝廷提出了什么；一份边塞简牍最适合说明特定关塞或官署留下了什么；一座城址、一枚钱币或一块残石能校验空间与物质条件；后出的史书、碑刻和宗教文献则各自保存编纂者、纪念者或社群的视角。问题决定材料的组合，而不是书名替代判断。

\begin{description}
  \sourceindexentry{史记}{秦汉转换、人物行动与战争叙事的重要入口。它的场面感不等于每一段言辞、数字或动机都可当作现场记录。}
  \sourceindexentry{汉书}{西汉帝纪、表、志与列传的基础材料；其以东汉形成的秩序回望西汉，尤其需要辨认正统与褒贬的组织方式。}
  \sourceindexentry{后汉书}{东汉政治、人物、制度与边疆的主干入口。范晔后出编纂，须与《后汉纪》、文书、遗存及不同地方传统交叉阅读。}
  \sourceindexentry{三国志}{汉末、曹魏与禅代问题的基础材料；不把魏、蜀、吴各自的叙事位置混作一个无立场的记录。}
  \sourceindexentry{资治通鉴}{连接年序、比照前后关系的编年参照，不替代更接近具体制度、地点或文书问题的证据。}
  \sourceindexentry{盐铁论}{观察财政、边防与民生如何被论辩的关键文本；不是前八一年会议的逐字速记，也不是全国意见调查。}
\end{description}

\section{文书、遗存与现代研究方向}

居延、肩水金关、悬泉置等简牍可使河西的军粮、邮传、出入关和吏员工作显形；其保存地点和机构性质也限制了它们对全帝国的代表性。钱范、钱币窖藏、城址、石经残石、碑刻、墓葬与佛道文本分别照出生产、流通、城市、教育、纪念和信仰的不同尺度。它们不应被降格为正史的插图，也不能由一件遗存推出全国人口、市场或信众的总量。

现代研究在本卷主要用于追问三类问题：第一，国家如何通过文书、运输和任用获取资源；第二，边地、地方社会与迁徙者如何改变政策的实际效果；第三，正史、出土材料和后出叙事之间哪里相合、哪里相抵。遇到兵数、疆界、族名、动机或仪式自愿性，读者应回到相应史书、文书或遗存的出处，比较它们各自看见了什么；叙事读来连贯，并不使这些问题自动获得确定答案。
